\section{Related Work}\label{sec:related}

Social bot detection has evolved from handcrafted feature engineering to graph-based representation learning and, more recently, reliability-aware modeling. Early methods relied on observable account, activity, network, and content signals, such as BotOrNot, online human--bot interaction modeling, Bot-Hunter, and FriendBot~\cite{davis2016botornot,varol2017online,bot_hunter2018,friendbot2020}, but these approaches remain vulnerable to feature manipulation and content imitation~\cite{cresci2017paradigm}. With the release of the TwiBot benchmarks~\cite{feng2021twibot20,twibot22benchmark}, graph-based methods increasingly exploited heterogeneous social relations as structural evidence for bot detection. Representative models include BotRGCN~\cite{feng2021botrgcn}, RGT~\cite{feng2022rgt}, and BIC~\cite{lei2023bic}, which improve account classification through relational propagation and text--graph interaction. Since pairwise edges may not fully capture coordinated behavior, recent work has further introduced higher-order relation modeling, for example by constructing hyperedges from neighborhood differences~\cite{shi2024hypergraphbot} or jointly modeling pairwise, hop-wise, and group-wise relations~\cite{gao2025higherorder}; however, noisy hypergraphs may also introduce task-irrelevant connections~\cite{cai2022hypergraphstructure}. More recent studies have begun to emphasize reliability in social bot detection. BECE estimates edge confidence to suppress unreliable graph evidence~\cite{qiao2025bece}, BotBR combines balanced feature fusion with reliability-enhanced graph learning~\cite{lin2025botbr}, and BotPercent calibrates model outputs for community-level bot-population estimation rather than node-level prediction correction~\cite{tan2023botpercent}. In contrast, our method uses a calibration-inspired local reliability estimate to guide selective higher-order correction for individual accounts.
